% !PRINTED 80
\weberpage{80}

%<<CONT
is the root of an Abelian equation whose group is isomorphic with the group $\mathfrak{B}$ reciprocal to $\mathfrak{A}$.

The problem that we now wish to pose and solve is the following:

\gesperrt{I. The modulus $m$ and the group $\mathfrak{A}$ are to be determined in all possible ways so that $\mathfrak{B}$ has an arbitrarily given system of invariants.}

Or, what is the same thing:

\gesperrt{All cyclotomic fields of given group are to be determined.}

In this we always make the assumption that $\mathfrak{A}$ is to be a primary divisor of $\mathfrak{N}$, because otherwise we would obtain one and the same field several times. The problem splits into two parts, namely:

\gesperrt{II. Which moduli $m$ are suited to generate a group $\mathfrak{B}$ of given invariants?}

\gesperrt{III. How does one find this group $\mathfrak{B}$ and the associated group $\mathfrak{A}$?}

In this investigation we must alter the notation somewhat as against the earlier, so that clarity is not lost. The exceptional position of the prime 2, which always had to be taken into account in the considerations hitherto, made a somewhat cumbersome notation necessary. In what follows this distinction is no longer necessary to the same degree, and therefore the notation can now be simplified.

We denote the indices of a number $a$ from the group $\mathfrak{A}$ by

\[
\alpha_1, \alpha_2, \dots \alpha_\mu \quad \text{(indices of $a$)}
\]

and the corresponding index moduli by

\[
c_1, c_2, \dots c_\mu \quad \text{(index moduli)},
\]

so that $\mu$ is equal to the number of primes dividing $m$, or, if $m$ is divisible by 8, greater by 1.

Let further

\[
\omega_1, \omega_2, \dots \omega_\mu
\]

be primitive roots of unity of the degrees $c_1, c_2, \dots c_\mu$. Now if the indices of a number $b$ in the group $\mathfrak{B}$ reciprocal to $\mathfrak{A}$ are

\[
\beta_1, \beta_2, \dots \beta_\mu \quad \text{(indices of $b$),}
\]
%CONT>>
